\documentclass[11pt,a4paper]{article}

% ---------- Kódování / jazyk / fonty (XeLaTeX) ----------
\usepackage{fontspec}
\usepackage{polyglossia}
\setdefaultlanguage{czech}
\setmainfont{DejaVu Serif}[Scale=0.92]
\setsansfont{DejaVu Sans}[Scale=0.92]
\setmonofont{Latin Modern Mono}[Scale=0.90]

% ---------- Uvozovky ----------
% \enquote{...} sází české „…“; autostyle navazuje na jazyk z polyglossia.
\usepackage[autostyle=true,czech=quotes]{csquotes}

% ---------- Geometrie ----------
\usepackage[a4paper,margin=1.8cm,headheight=14pt]{geometry}

% ---------- Barevná paleta (shodná s docx variantou) ----------
\usepackage[table,dvipsnames]{xcolor}
\definecolor{navy}{HTML}{1F3A5F}
\definecolor{blue}{HTML}{2E75B6}
\definecolor{light}{HTML}{D5E8F0}
\definecolor{grey}{HTML}{F2F2F2}
\definecolor{amber}{HTML}{FFF2CC}
\definecolor{green}{HTML}{E2EFDA}
\definecolor{red}{HTML}{FBE4E4}
\definecolor{rule}{HTML}{CCCCCC}

% ---------- Nadpisy ----------
\usepackage{titlesec}
\titleformat{\section}{\sffamily\Large\bfseries\color{navy}}{\thesection.}{0.6em}{}
\titleformat{\subsection}{\sffamily\large\bfseries\color{blue}}{\thesubsection}{0.6em}{}
\titleformat{\subsubsection}{\sffamily\normalsize\bfseries\color{black!70}}{\thesubsubsection}{0.5em}{}
\titlespacing*{\section}{0pt}{14pt}{8pt}

% ---------- Hlavička / patička ----------
\usepackage{fancyhdr}
\usepackage{lastpage}
\pagestyle{fancy}
\fancyhf{}
\renewcommand{\headrulewidth}{0.4pt}
\renewcommand{\footrulewidth}{0pt}
\fancyhead[L]{\sffamily\footnotesize\color{black!45}AI-asistovaný vývoj — Pravidla pro projekty}
\fancyfoot[C]{\sffamily\footnotesize\color{black!45}Interní — \versionline \quad Strana \thepage{} z \pageref{LastPage}}

% ---------- Tabulky ----------
\usepackage{tabularx}
\usepackage{booktabs}
\usepackage{array}
\usepackage{ragged2e}
\renewcommand{\arraystretch}{1.35}
\newcolumntype{L}[1]{>{\RaggedRight\arraybackslash}p{#1}}
% buňka blast-radius: \br{barva}{text}
\newcommand{\br}[2]{\cellcolor{#1}#2}

% ---------- Callout boxy ----------
\usepackage[most]{tcolorbox}
\newtcolorbox{callout}[2][light]{
  enhanced,
  colback=#1,colframe=blue,
  boxrule=0pt,leftrule=3pt,
  arc=1pt,left=8pt,right=8pt,top=6pt,bottom=6pt,
  breakable,
  coltitle=white,
  colbacktitle=blue,
  fonttitle=\sffamily\bfseries,
  attach boxed title to top left={xshift=3pt,yshift=-2pt},
  boxed title style={colframe=blue,arc=1pt,boxrule=0pt},
  title={#2}
}
\newtcolorbox{calloutplain}[1][grey]{
  colback=#1,colframe=blue,
  boxrule=0pt,leftrule=3pt,
  arc=1pt,left=8pt,right=8pt,top=6pt,bottom=6pt,breakable
}

% ---------- Část-banner (ČÁST A / ČÁST B) ----------
\newtcolorbox{partbanner}[2][navy]{
  colback=#1,colframe=#1,
  boxrule=0pt,arc=2pt,left=10pt,right=10pt,top=8pt,bottom=8pt,
  fontupper=\sffamily\Large\bfseries\color{white},
  before skip=14pt,after skip=10pt
}

% ---------- Seznamy ----------
\usepackage{enumitem}
\setlist[itemize]{leftmargin=1.4em,itemsep=2pt,topsep=3pt}
\setlist[enumerate]{leftmargin=2.2em,itemsep=3pt,topsep=3pt}

% ---------- Checkbox ----------
\usepackage{amssymb}
\newcommand{\cb}{$\square$\hspace{0.4em}}

% ---------- Kód / listing pro ADR ----------
\usepackage{fancyvrb}
\usepackage{fvextra}

% ---------- Odkazy ----------
\usepackage[hidelinks]{hyperref}
\hypersetup{
  pdftitle={AI-asistovaný vývoj — Pravidla pro projekty},
  pdfsubject={Interní směrnice + best practices},
  bookmarksnumbered=true
}

% ---------- Pomocná makra ----------
\newcommand{\code}[1]{{\ttfamily #1}}

% ---------- Verze dokumentu (vstřikuje CI) ----------
% CI vygeneruje version.tex s \def\docversion / \doccommit / \docdate.
% Lokální build bez version.tex spadne na fallback "draft".
\providecommand{\docversion}{draft}
\providecommand{\doccommit}{local}
\providecommand{\docdate}{\today}
\IfFileExists{version.tex}{\input{version.tex}}{}
% Patičkový a titulkový řetězec verze
\newcommand{\versionline}{v\docversion\ (\doccommit) — \docdate}

% ---------- Sazba: tolerantnější lámání řádků ----------
\tolerance=1200
\emergencystretch=2em
\hyphenpenalty=50
\hbadness=2000
